\subsection*{導入：ソフトウェア構成管理が必要な理由}
\addcontentsline{toc}{subsection}{導入：ソフトウェア構成管理が必要な理由}
\term{ソフトウェア構成管理}{Software Configuration Management}は、ソフトウェアライフサイクル全体を通じて構成管理を適用し、構成品目の完全性と正確性を保つ活動である。構成管理は、構成品目の機能的・物理的特性を識別して文書化し、それらへの変更を統制し、変更処理と実施状況を記録・報告し、指定要求への適合を検証するための技術的・管理的な規律である。

SCM は、プロジェクト管理、開発、保守、品質保証、顧客、最終利用者を支える。開発方法がウォーターフォール型であっても、反復・増分型であっても、管理対象を明確にし、変更と版を管理し、証跡を残す必要がある点は変わらない。

SCM はソフトウェア品質保証と密接に関係する。品質保証は、成果物とプロセスが要求に合っていることを示す活動である。SCM は、構成品目の識別、変更統制、状態記録、監査を通じて、その証拠を残す。

\begin{notebox}{重要}SCM の目的は、変更を禁止することではない。変更を安全に行い、何が、誰によって、いつ、なぜ、どの版に対して変わったのかを後から説明できるようにすることである。\end{notebox}
\par\smallskip
\begin{center}
\centering
\IfFileExists{assets/generated_figures/ch08/fig-ch08-02.png}{\includegraphics[width=0.96\linewidth]{assets/generated_figures/ch08/fig-ch08-02.png}}{\fbox{\parbox[c][48mm][c]{0.9\linewidth}{\centering 画像生成待ち\\SCM が支える活動の関係図}}}
\par\smallskip
\refstepcounter{figure}\label{fig:ch08-02}図 \thefigure: SCM が支える活動の関係図
\end{center}
図 \thefigure は、SCM が支える活動の関係図を視覚的に整理したものである。
\par\smallskip
